\emph{Setting.} The derived metric is $d\tau^{2} - \sum_{j}\tau^{2p_{j}}(dx^{j})^{2}$ with $\sum_{j}p_{j} = 1$, so $\sqrt{|g|} = \tau$ and the expansion rate (the logarithmic derivative of the volume) is $\theta = (\sum_{j}p_{j})/\tau = 1/\tau$, the Hubble length $1/\theta = \tau$. The momentum $k$ lies along the principal axis $i$ with $p_{i} < 1$, $q = 1 - p_{i} > 0$ and $s = k\tau^{q}/q$, as in Theorem \claimref{synthesis.kasner-approach-gives-horizon-exit-and-index-three}. By Theorem \claimref{synthesis.kasner-approach-coupled-modes-reach-measured-tilt-only-by-a-fitted-curvature-coupling}(ii) every mode treated in the store's Kasner results (the minimally coupled scalar, curvature-coupled scalars at disk points, the gravitational field's off-diagonal modes, and, by the same reduction, scalars weighted by the hidden scale) is $s^{-w/(2q)}Z(s)$ with $Z$ obeying Bessel's equation of some order $\nu \ge 0$; the order is fixed by the field and the point, and nothing below depends on its value except where stated.

\emph{(i)} The physical wavelength along the axis is $\tau^{p_{i}}/k$ (comoving $1/k$ times the scale $\tau^{p_{i}}$), so its ratio to the Hubble length is $\tau^{p_{i}}/(k\tau) = 1/(k\tau^{q}) = 1/(qs)$. Since $q > 0$, $d(k\tau^{q})/d\tau = kq\tau^{q - 1} > 0$, so the ratio is strictly decreasing in $\tau$: it falls only as $\tau$ falls. A mode is outside the Hubble length when the ratio exceeds one, so modes leave only on the branch along which $\tau$ decreases toward the wall, and along a branch on which $\tau$ grows from the wall every mode begins outside (the ratio tends to infinity as $\tau \to 0$) and enters. The order $\nu$ does not appear. The check samples seven exponents ($p_{i} \in \{-\tfrac13, \tfrac23, \tfrac13, 0, \tfrac12, -\tfrac27, \tfrac67\}$) and five values of $\tau$ chosen as perfect powers so that $\tau^{q}$ is rational, and confirms the ratio $1/(k\tau^{q})$ and its monotonicity exactly. Theorem \claimref{synthesis.kasner-approach-gives-horizon-exit-and-index-three} closes with this remark for order $0$; the computation shows that it holds for every order, which is all (i) adds.

\emph{(ii)} Put $Z = s^{-1/2}u$. Then $Z' = s^{-1/2}u' - \tfrac12 s^{-3/2}u$ and $Z'' = s^{-1/2}u'' - s^{-3/2}u' + \tfrac34 s^{-5/2}u$, so
\[
Z'' + \frac{Z'}{s} + \Bigl(1 - \frac{\nu^{2}}{s^{2}}\Bigr)Z = s^{-1/2}\Bigl(u'' + \Bigl(1 + \frac{\tfrac14 - \nu^{2}}{s^{2}}\Bigr)u\Bigr),
\]
the Liouville form (Fact \claimref{synthesis.imported-bessel-reduction-and-small-argument-hankel-forms} for the Bessel equation itself); the check verifies the identity on monomials $u = s^{m}$ for five orders and five exponents, with rational exponents, and a control with the weight $s^{+1/2}$ fails. Write $\omega^{2} = 1 + a/s^{2}$, $a = \tfrac14 - \nu^{2}$. Adiabaticity is taken in the Liouville form in $s$: an adiabatic (positive-frequency WKB) state is defined where $\omega^{2} > 0$ and $|\omega'/\omega^{2}| \ll 1$. For $0 \le \nu < \tfrac12$, $0 < a \le \tfrac14$, $\omega^{2} > 0$, $d\omega^{2}/ds = -2a/s^{3}$, so $\omega' = -a/(s^{3}\omega)$ and
\[
\Bigl(\frac{\omega'}{\omega^{2}}\Bigr)^{2} = \frac{a^{2}}{s^{6}\omega^{6}} = \frac{a^{2}}{(s^{2} + a)^{3}} ,
\]
which is $1/a \ge 4$ at $s = 0$ and strictly decreasing in $s$; it is at least $1$ exactly for $s^{2} \le a^{2/3} - a$ (positive since $a < 1$). The check confirms the value at the wall for four values of $a$, the equality $a^{2}/(s^{2} + a)^{3} = 1$ at $s^{2} = a^{2/3} - a$ for three cubes $a$, and the monotonicity on a grid. For $\nu > \tfrac12$, $a < 0$ and $\omega^{2} < 0$ for $s^{2} < -a = \nu^{2} - \tfrac14$: no oscillation, so no positive-frequency condition, near the wall (checked at three orders). For $\nu = \tfrac12$, $a = 0$ and $u'' + u = 0$ exactly; the positive-frequency solution $u = e^{-I_4 s}/\sqrt{2}$ up to normalisation (the field theory's unit $I_4$; numbers are read in $\mathrm{span}\{1, I_4\}$) gives $|Z|^{2} \propto 1/s$, so with $|\chi_{k}|^{2} \propto \tau^{-w}|Z|^{2}$ and $s \propto k$ the dimensionless spectrum is $k^{3}/k = k^{2}$, index $2 = 3 - 2\nu$. Therefore on a branch expanding from the wall, where by (i) every mode starts at $s \to 0$, no mode with $\nu \ne \tfrac12$ has its state fixed by an adiabatic condition in $s$ at the wall (in the variable $\ln s$ the equation for $Z$ is $Z_{\ln s \ln s} + (s^{2} - \nu^{2})Z = 0$, whose frequency squared $s^{2} - \nu^{2}$ is negative or tends to zero at the wall, with the same conclusion): its two-point function at the wall is initial data, carried forward by the mode equation. The stored spectra of the Kasner results are computed in the opposite direction, from an adiabatic state inside ($s \gg 1$) toward the wall.

\emph{(iii)} For the coframe $e^{0} = d\tau$, $e^{i} = \tau^{p_{i}}dx^{i}$ the connection forms are $\omega^{i0} = (p_{i}/\tau)e^{i}$, $\omega^{ij} = 0$, and Cartan's second structure equation gives the curvature two-forms $R^{i0} = p_{i}(p_{i} - 1)\tau^{-2}e^{0}\wedge e^{i}$ and $R^{ij} = -p_{i}p_{j}\tau^{-2}e^{i}\wedge e^{j}$, the formulas of Theorem \claimref{dark-energy.frozen-point-datum-kasner-wall}(i), which follow from the power-law coframe alone with no field equation, for any number of directions; the check recomputes them from the connection forms, $d\omega^{i0}$ giving $p_{i}(p_{i} - 1)\tau^{p_{i} - 2}d\tau\wedge dx^{i}$ and $\omega^{i0}\wedge\omega^{0j}$ the coefficient $p_{i}p_{j}\tau^{p_{i} + p_{j} - 2}$, divided by the frame factors. The curvature, as a map of bivectors, is diagonal on $e^{0}\wedge e^{i}$ and $e^{i}\wedge e^{j}$, so $R_{abcd}R^{abcd} = 4\sum$ of the squares of these coefficients:
\[
K\tau^{4} = 4\sum_{i}p_{i}^{2}(p_{i} - 1)^{2} + 4\sum_{i<j}p_{i}^{2}p_{j}^{2} .
\]
On the four-dimensional vacuum circle this equals $-16p_{1}p_{2}p_{3}$, the stored Kretschmann scalar (checked at four values of the BKL parameter; $64/27$ at the Schwarzschild exponents $(-\tfrac13, \tfrac23, \tfrac23)$); a control with the sign of the pair term reversed fails. $K\tau^{4}$ is a sum of squares, so it vanishes exactly when every $p_{i}(p_{i} - 1) = 0$ and every $p_{i}p_{j} = 0$ ($i \ne j$): each exponent is $0$ or $1$ and at most one is nonzero, so with $\sum_{j}p_{j} = 1$ exactly one is $1$, a flat Milne point. Every other Kasner wall has $K \propto \tau^{-4} \to \infty$, a curvature singularity in every chart since $K$ is a scalar; the check confirms $K > 0$ at the isotropic disk point $(\tfrac13, \tfrac13, \tfrac13)$, at $(\tfrac12, \tfrac12, \tfrac12, -\tfrac12)$, at the generic five-dimensional vacuum point $(\tfrac15, \tfrac25, \tfrac45, -\tfrac25)$ and at the four-dimensional vacuum point $(-\tfrac27, \tfrac37, \tfrac67)$ embedded with a static hidden direction, and $K = 0$ at $(1, 0, 0)$ and $(0, 0, 0, 1)$. At $(0, 0, 0, 1)$ the five-dimensional metric $d\tau^{2} - (dx)^{2} - \tau^{2}dy^{2}$ is flat, the hidden scale $b = \tau$ tends to zero, and the reduction of Theorem \claimref{synthesis.scalar-kasner-wall-is-reduced-five-dimensional-vacuum-wall}, $p_{i} = (P_{i} + P_{4}/2)/(1 + P_{4}/2)$, sends it to $(\tfrac13, \tfrac13, \tfrac13)$ (checked), where the four-dimensional $K$ diverges.

\emph{Reading (argued).} Two hypotheses frame the reading: the dating's placement of the descent at the origin (\claimref{standard-model.dating-descent}), and that the descended phase expands monotonically from the descent, with no recollapse and bounce after it (the record supports expansion back to nucleosynthesis; without this hypothesis a contracting stage after the descent is not excluded by (i)). Under them a Kasner stage inside the descended phase that begins at the descent is a branch expanding from a wall: by (i) no mode exits on it, and by (ii) the state of its modes at the wall is initial data, which for the descended phase is the descent's deposited state, the open datum of Theorem \claimref{synthesis.finite-range-descent-deposit-cannot-seed-cmb}. Such a stage adds no source of its own. An exited Kasner spectrum therefore needs (P1) a metric stage contracting toward a wall, which with the descent at the origin lies before the descent, and (P2) a rule carrying the exited state from that contracting stage into the descended phase. If the approach runs to the wall, by (iii) the wall is a curvature singularity at every stored Kasner point except the five-dimensional Milne end, and, by the argument of Theorem \claimref{dark-energy.frozen-point-datum-kasner-wall}(i) (a diverging curvature scalar), the frozen-point datum extends regularly across it in no chart. If the descent ends the approach at finite $\tau$, after the modes have left, the rule needed is a matching across the descent itself. In either case, under Einstein's equations a turn from contraction to expansion needs a singularity or a violation of the null energy condition (standard), so (P2) is needed whatever ends the approach; the store holds no such rule. The dating's reading (iii) takes pre-descent physics to be the pre-metric sector, the polynomial action and the $\theta$-remnants alone, which denies (P1); that reading has the status of a conjecture, and no proved result denies (P1): gravity's polynomial action is defined across the singular locus (\claimref{dark-energy.gravity-pre-metric-forces-degenerate-exactly}(2)), so a metric field before the descent is not excluded by the action. A pre-descent Kasner stage is thus excluded by the dating's reading and admitted only with (P1) and (P2) as hypotheses; the tension between them and the pre-metric reading is recorded here, not resolved. Granting (P1) and (P2) does not supply the tilt: at tree level every curvature coupling the series fixes leaves the exited spectrum blue at the Kasner points the store fixes (\claimref{synthesis.tree-level-curvature-couplings-stay-blue-at-the-isotropic-kasner-point}), so the coupling that matches the measured tilt remains a free parameter. At the five-dimensional Milne end the wall is the collapse of the hidden circle in flat five-dimensional spacetime, the compactified Milne geometry in which the ekpyrotic and cyclic models place their brane collision; the continuation of perturbations through it is not settled in that literature either (Khoury, Ovrut, Steinhardt and Turok 2001; Tolley, Turok and Steinhardt 2004; cited from memory).

\emph{Level.} (i)-(iii) proved by the displayed computations, which the check confirms exactly at the sampled points; the general statements follow from the closed forms. The reading is argued. Transport-absent.
